\section{Cap estimates in a first-order cone chart}\label{weak:section}

Fix \(\ell\ge 1\). The comparison uses a single chart through the
splices and both retained buffers. Write
\[
 A_\ell=\{y\in\mathbb R^4:e^{-\ell-2}<|y|<e^{\ell+2}\},
 \qquad t=\log|y| .
\]
An oriented smooth embedding \(\Phi:A_\ell\to M\) puts the two
splices \(T_j\) at \(0<|t|<1\) and the retained buffers at
\(1<|t|<1+\ell\). For a constant \(\rho>0\), use the metric
\(m=\rho^{-2}\Phi^*g=\delta_{\rm Euc}+v\) in these Cartesian
coordinates. Set
\[
 a=\|v\|_{L^\infty(A_\ell)},\qquad
 d=\|Dv\|_{L^4(A_\ell,dy)},\qquad
 K=\|\operatorname{Rm}_m\|_{L^2(T_1\cup T_2,m)} .
\]
The first two norms use Euclidean components, and the third uses
the four-index curvature norm and volume of \(m\). The last norm
equals the corresponding norm of \(g\) on the embedded splices,
by four-dimensional scaling. On the product collar use
\(h=|y|^{-2}m\), extending the positive conformal multiplier to the
source outside the region used below. On each splice put
\(n=\delta_{\rm Euc}+\theta(|t|)v\), with the cutoff of
Lemma~\ref{lem:buffered-caps}, and attach the same model cap to
\(|y|^{-2}n\). This is the exact product near the middle face and the source representative near the outer splice face.

\begin{lemma}[First-order cap comparison]\label{weak:splice}
Suppose \(a\le 1/4\), \(Y([g])>0\), and \(C_\ell(a+d)\le 1/2\)
for a sufficiently large constant \(C_\ell\). Write \(P_j\)
for the raw sides and \(\widehat g_j\) for the capped metrics.
For either chirality let \(q_j=E_\pm(g;T_j)\) and
\(q=q_1+q_2\). There is an error
\[
 \Delta=C\bigl(aK+d^2+d+a\bigr)
\]
such that
\begin{align}
 \left(\sum_j E_\pm(\widehat g_j;T_j)\right)^{1/2}
     &\le \sqrt q+\Delta,\\
 \sum_j E_\pm(\widehat g_j)
     &\le E_\pm(g)+2\sqrt q\,\Delta+\Delta^2,\\
 E_\pm(\widehat g_j)
     &\le E_\pm(g;P_j)+2\sqrt{q_j}\,\Delta+\Delta^2 .
 \label{weak:side-cost}
\end{align}
The Yamabe constants satisfy
\begin{equation}\label{weak:yamabe}
 Y([\widehat g_j])\ge
 \min\{Y([g]),(1-C_\ell(a+d))Y_{\rm rd}\}
 - C_I\ell^{-3/2}.
\end{equation}
The constants \(C,C_I\) depend only on the fixed profiles and
model, and \(C_\ell\) may also depend on \(\ell\).
\end{lemma}

\begin{proof}
For an all-lowered curvature tensor in Cartesian coordinates,
\[
 \operatorname{Rm}_{\rm low}(p)=L(D^2p)+Q(p)[Dp,Dp]
\]
with \(L\) constant-coefficient and \(Q\) smooth on positive
metric matrices. The second derivatives arise by differentiating
Christoffel symbols of the first kind, which are linear in
\(Dp\) with constant coefficients. Differentiation of the inverse
metric and the products of Christoffel symbols contribute to \(Q\).
The principal terms in
\(\operatorname{Rm}_{\rm low}(n)-\theta\operatorname{Rm}_{\rm low}(m)\)
cancel except for the product-rule terms containing
\(D\theta\,Dv\) and \(D^2\theta\,v\). The Cartesian derivatives
of \(\theta\) are bounded on the unit splices. Thus the difference
is bounded by \(C(|Dv|^2+|Dv|+a)\).

Let \(\mathcal P_p\) take an all-lowered curvature tensor to its
chiral Weyl matrix in the oriented \(p^{-1/2}\)-orthonormal frame,
multiplied by \((\det p)^{1/4}\). This linear map depends smoothly
on \(p\), so \(\|\mathcal P_n-\mathcal P_m\|\le Ca\).
The square of the Euclidean \(L^2\) norm of
\(F(p)=\mathcal P_p\operatorname{Rm}_{\rm low}(p)\) is the chiral
energy. Uniform metric comparison and H\"older yield
\[
 \|F(n)-\theta F(m)\|_{L^2(T_1\cup T_2,dy)}
 \le C\bigl(d^2+d+a+aK\bigr).
\]
Indeed the curvature difference is acted on by \(\mathcal P_n\),
and the error from changing the projection is
\(\theta(\mathcal P_n-\mathcal P_m)\operatorname{Rm}_{\rm low}(m)\).
This proves the square-root estimate on the union of the splices,
and on each one with its own incoming energy. Chiral energy is
conformally invariant. The net cost is the difference of the new
and incoming splice energies, as in \eqref{eq:cut-splice-excess}.

For a localized test on the cap side, extend the perturbation to the
model \(b=b_\infty\) of \eqref{eq:cut-model}, obtaining \(k=b+e\).
Use the actual capped metric through the buffer, then cut
\(e\) off on \(1+\ell<|t|<2+\ell\).
The perturbation is zero on the exact cap. In product coordinates
\(b=|y|^{-2}\delta_{\rm Euc}\), so the metric, connection and
cutoff comparisons give
\[
 \|e\|_{L^\infty(b)}\le a,\qquad
 D_e:=\|\nabla_b e\|_{L^4(b)}\le C\bigl(d+(1+\ell)^{1/4}a\bigr).
\]
In particular the product volume density is \(|y|^{-4}dy\)
and \(|\nabla_b e|_b\le C(|y||Dv|+|v|)\) on the support of \(e\).
Put \(P^{ij}=(dV_k/dV_b)k^{ij}\) and
\(A^p{}_{ij}=(\nabla_k-\nabla_b)^p{}_{ij}\).
Then \(|P-b^{-1}|\le C|e|\) and
\(|\nabla_bP|+|A|\le C|\nabla_b e|\) in model norms.
Contracting the Ricci difference with \(P\) gives a divergence
with vector \(V^p=P^{ij}A^p{}_{ij}-P^{ip}A^q{}_{iq}\).
The derivatives of \(P\) and the connection products are
quadratic in \(\nabla_b e\), and the change in contraction of
the background Ricci tensor is \(O(|e|)\). Hence
\[
 (dV_k/dV_b)R_k-R_b=\operatorname{div}_b V+H,\qquad
 |V|\le C|\nabla_b e|,\qquad |H|\le C(|e|+|\nabla_b e|^2).
\]
Use \(Q_p(w)=\int(6|dw|_p^2+R_p w^2)dV_p\).
Integrating the divergence against \(w^2\) for a compactly
supported smooth test gives
\[
 |Q_k(w)-Q_b(w)|\le C\bigl(a\|w\|_{H^1(b)}^2+
 D_e\|w\|_{4,b}\|dw\|_{2,b}+D_e^2\|w\|_{4,b}^2\bigr).
\]
The positive scalar lower bound and the round Yamabe inequality
on the model bound this error by \(C_\ell(a+d)Q_b(w)\).
The denominator changes by a factor at most \(1+Ca\), giving
\(Q_k(w)\ge (1-C_\ell(a+d))Y_{\rm rd}\|w\|_{4,k}^2\)
after increasing \(C_\ell\).

Apply \eqref{eq:cut-ims} with the source cutoff supported away from
the splice and the cap cutoff vanishing near the far edge of the
buffer. Their tests extend by zero to the conformal source and to
\(k\), respectively. Uniform comparison to the product on the
buffer bounds the sum of squared cutoff gradients by
\(C\ell^{-2}\) on volume \(C\ell\). Its \(L^2\) norm is at most
\(C\ell^{-3/2}\). The triangle inequality for the squares of the
localized functions then gives \eqref{weak:yamabe}.
\end{proof}

\begin{corollary}[Measured balanced sides]\label{weak:balanced}
Let \(g_i\) be metrics on a fixed oriented smooth homotopy
four-sphere \(M\), with \(Y([g_i])\ge y_0>0\) and
\(E_+(g_i)\le 8\pi^2-\delta\) for some \(\delta>0\).
Choose \(\ell\) large enough that
\(C_I\ell^{-3/2}<\tfrac12\min\{y_0,Y_{\rm rd}/2\}\).
Suppose the charts above are supplied at this fixed length, with
\(a_i+d_i\to0\), \(K_i\) bounded, and measurable sets
\(D_{i,j}\subset P_{i,j}\) satisfying
\(E_+(g_i;D_{i,j})\ge\beta>0\) for both sides.
Then both caps have positive Yamabe constant for large \(i\), and
each has positive-chiral energy at most \(E_+(g_i)-\beta+o(1)\).
If moreover \(E_+(g_i)\to I^+(M)\), the infimum over positive
Yamabe conformal classes on \(M\), then \(M\) is diffeomorphic
to \(S^4\).
\end{corollary}

\begin{proof}
Equation~\eqref{weak:yamabe} gives positivity. Each raw side omits
at least \(\beta\) in the other side, and each surcharge in
\eqref{weak:side-cost} tends to zero. Thus \(E_+(g_i)\) plus the
sum of the two surcharges is strictly below \(8\pi^2\) for large
\(i\). If \(M\) is nonstandard, \eqref{eq:cut-discarded} in
Proposition~\ref{prop:cut-energy} now gives
\[
 \min_j E_+(g_i;P_{i,j})\le E_+(g_i)-I^+(M)+o(1).
\]
This contradicts the fixed \(\beta\) margin along a minimizing
sequence.
\end{proof}
